\documentclass[journal]{IEEEtran}

% *** PACKAGES ***
\usepackage{amsmath}
\usepackage{amssymb}
\usepackage{graphicx}
\usepackage{booktabs}
\usepackage{multirow}
\usepackage{array}
\usepackage{url}
\usepackage{cite}
\usepackage{algorithmic}
\usepackage{algorithm}
\usepackage{xcolor}
\usepackage{subfig}

\hyphenation{op-tical net-works semi-conduc-tor water-mark-ing}
\newcommand{\safeincludegraphics}[2]{%
  \IfFileExists{#2}{%
    \includegraphics[width=#1]{#2}%
  }{%
    \IfFileExists{figures/#2}{%
      \includegraphics[width=#1]{figures/#2}%
    }{%
      \IfFileExists{../figures/#2}{%
        \includegraphics[width=#1]{../figures/#2}%
      }{%
        \fbox{\parbox[c][4cm][c]{#1}{\centering Missing image file:\\\texttt{#2}}}%
      }%
    }%
  }%
}

\begin{document}

% ============================================================
% TITLE
% ============================================================
\title{CatNMF-Net: Image Watermarking Using Catalan Transform and Non-Negative Matrix Factorization}

\author{Author~One,~\IEEEmembership{Member,~IEEE,}
        Author~Two,~\IEEEmembership{Member,~IEEE,}
        and~Author~Three,~\IEEEmembership{Senior~Member,~IEEE}%
\thanks{Author One is with the Department of Computer Science and
Engineering, XYZ University, City, Country (e-mail: author1@xyz.edu).}%
\thanks{Author Two and Author Three are with the Department of
Electrical Engineering, ABC University, City, Country.}%
\thanks{Manuscript received Month DD, YYYY; revised Month DD, YYYY.}}

\markboth{IEEE TRANSACTIONS ON INFORMATION FORENSICS AND SECURITY,~VOL.~XX,
NO.~XX, MONTH~YEAR}%
{Author \MakeLowercase{\textit{et al.}}: CatNMF-Net: Deep Learning Based
Image Watermarking Using Catalan Transform and NMF}

\maketitle

% ============================================================
% ABSTRACT
% ============================================================
\begin{abstract}
This paper presents \textbf{CatNMF-Net}, a novel deep learning framework
for robust and imperceptible digital image watermarking that integrates
three complementary mathematical foundations: the Haar Discrete Wavelet
Transform (DWT), a QR-orthogonalized Catalan Transform, and Non-Negative
Matrix Factorization (NMF). The proposed architecture decomposes a
grayscale cover image into wavelet subbands, applies a block-wise Catalan
Transform to the low-frequency approximation subband (LL), and embeds a
binary watermark through a residual attention-based encoder network.
Security is reinforced by an NMF decomposition layer that produces a
private ownership key from the watermarked LL subband, enabling
cryptographically motivated ownership verification. A curriculum training
strategy is employed to sequentially prioritize watermark fidelity before
image quality optimization, preventing the well-known embedding-collapse
phenomenon in deep watermarking systems. Extensive experiments on the
DIV2K high-resolution benchmark dataset demonstrate that CatNMF-Net
achieves a Normalized Cross-Correlation (NCC) of \textbf{0.9998}, a Bit
Error Rate (BER) of \textbf{0.0132}, and a Peak Signal-to-Noise Ratio
(PSNR) of \textbf{25.48~dB}. The system successfully withstands 12 out
of 20 diverse attack types including JPEG compression, Gaussian noise,
salt-and-pepper noise, and histogram equalization. The NMF-based ownership
verification achieves an NCC of 0.9164 for genuine keys versus 0.3608
for forged keys, confirming strong security guarantees. An ablation study
validates the necessity of the Catalan Transform, showing a 35.8\%
degradation in NCC when the transform is omitted. The complete model
requires only 608,244 trainable parameters, making it practically
deployable in resource-constrained environments.
\end{abstract}

\begin{IEEEkeywords}
Image watermarking, Catalan transform, non-negative matrix factorization,
discrete wavelet transform, deep learning, attention mechanism, robustness,
curriculum training, ownership verification, DIV2K.
\end{IEEEkeywords}

\IEEEpeerreviewmaketitle

% ============================================================
% SECTION I: INTRODUCTION
% ============================================================
\section{Introduction}

\IEEEPARstart{D}{igital} watermarking is a cornerstone technology in
multimedia security, providing mechanisms for copyright protection,
content authentication, and provenance tracking. A watermarking system
must satisfy four primary requirements that are fundamentally in tension:
\emph{imperceptibility} (the watermark must be invisible to human
observers), \emph{robustness} (the watermark must survive common signal
processing attacks), \emph{capacity} (sufficient bits must be embedded),
and \emph{security} (unauthorized parties must not be able to detect,
remove, or forge the watermark) \cite{cox1997secure}.

Classical frequency-domain methods exploit the perceptual properties of
transform coefficients. Spread-spectrum techniques embed watermarks in
the DCT domain \cite{cox1997secure}, while wavelet-domain methods
leverage the multi-resolution analysis properties of the DWT
\cite{barni1998dct}. Hybrid approaches combining DWT, DCT, and Singular
Value Decomposition (SVD) have achieved state-of-the-art imperceptibility
with PSNR values exceeding 40~dB \cite{makbol2013robust,parah2016hiding}.
Non-Negative Matrix Factorization (NMF) has emerged as a powerful
alternative to SVD, offering non-negativity constraints that align with
the physical interpretation of pixel intensities \cite{hamidi2018hybrid}.

The advent of deep learning has fundamentally transformed watermarking
research. Encoder-decoder convolutional architectures such as HiDDeN
\cite{zhu2018hidden} demonstrated that end-to-end differentiable systems
can simultaneously optimize embedding quality and extraction accuracy.
Subsequent works have incorporated adversarial training, attention
mechanisms, and transformer architectures to improve robustness against
geometric distortions and advanced attacks \cite{zhang2022robust}.
However, deep learning approaches often sacrifice interpretability and
security guarantees that frequency-domain methods naturally provide.

\subsection{Motivation and Contributions}

This paper bridges the gap between classical mathematical transforms and
modern deep learning by proposing CatNMF-Net, which leverages:

\begin{enumerate}
  \item \textbf{Catalan Transform}: A novel orthogonal transform derived
  from the Catalan number sequence, applied block-wise to the DWT LL
  subband. The QR-orthogonalized Catalan matrix concentrates energy
  similarly to the DCT while providing a unique, non-standard basis that
  offers an additional layer of security through transform obscurity.

  \item \textbf{NMF Security Layer}: A factorization-based ownership
  verification mechanism where $\mathbf{V} \approx \mathbf{W} \times
  \mathbf{H}$, with $\mathbf{W}$ serving as a private ownership key. The
  non-negativity constraint ensures physical interpretability and
  resistance to random key forgery.

  \item \textbf{Curriculum Training}: A staged loss weighting strategy
  that prevents the embedding-collapse phenomenon observed in prior
  deep watermarking systems, where the model learns to embed nothing
  (achieving perfect image quality at the cost of zero watermark
  fidelity).

  \item \textbf{Spatial Attention Embedding}: An adaptive mechanism that
  preferentially embeds the watermark in textured image regions, which
  are less sensitive to human visual perception, maximizing imperceptibility
  for a given embedding strength.
\end{enumerate}

The main contributions of this paper are:
\begin{itemize}
  \item We propose the first watermarking system to combine the Catalan
  Transform with deep learning, introducing a novel frequency-domain
  basis for embedding.
  \item We design a curriculum training protocol with noise-layer
  scheduling that resolves the embedding-collapse failure mode
  common in attention-based deep watermarking.
  \item We integrate NMF-based ownership verification as a post-embedding
  security layer, achieving an owner-vs-fake NCC discrimination ratio
  of $0.9164$ vs. $0.3608$.
  \item We conduct comprehensive robustness evaluation across 20 attack
  types and an ablation study validating each architectural component.
\end{itemize}

The remainder of this paper is organized as follows. Section II reviews
classical and deep learning watermarking methods and identifies open
gaps addressed by this study. Section III presents the complete CatNMF-Net
system, including transform design, network modules, and security
mechanisms. Section IV details optimization, curriculum scheduling, and
convergence stabilization. Section V defines dataset, protocol, and
evaluation criteria. Section VI reports quantitative and qualitative
results, including attack-wise and ablation analyses. Section VII
discusses practical deployment, ethical boundaries, and future research.
Section VIII concludes the paper.

% ============================================================
% SECTION II: RELATED WORK
% ============================================================
\section{Related Work}

\subsection{Frequency-Domain Watermarking}

The seminal work of Cox \textit{et al.} \cite{cox1997secure} established
DCT spread-spectrum watermarking as a robust baseline, achieving PSNR
above 35~dB with NCC approximately 0.75 under common attacks. Barni
\textit{et al.} \cite{barni1998dct} extended this to the DWT domain,
exploiting the perceptual masking properties of wavelet coefficients to
improve imperceptibility to $>$38~dB. The DWT+SVD paradigm introduced
by Makbol \textit{et al.} \cite{makbol2013robust} further improved NCC
to approximately 0.88 by embedding in the singular values of wavelet
subbands, achieving PSNR exceeding 40~dB.

Parah \textit{et al.} \cite{parah2016hiding} proposed a triple-transform
approach combining DWT, DCT, and SVD with inter-block embedding, reporting
PSNR $>$41~dB and NCC $\approx$0.90. The integration of NMF into
watermarking was pioneered by Hamidi \textit{et al.}
\cite{hamidi2018hybrid}, who demonstrated that NMF's non-negativity
constraints provide more stable embedding across diverse image content
compared to SVD, achieving NCC $\approx$0.92. Ernawan \textit{et al.}
\cite{ernawan2020robust} proposed a chaotic sequence-based coefficient
selection (CSBC) within the DWT+DCT framework, reporting NCC $\approx$0.93
on the CIPR dataset.

\subsection{Deep Learning Watermarking}

HiDDeN \cite{zhu2018hidden} established the encoder-noise-decoder
paradigm for end-to-end differentiable watermarking. The noise layer,
simulating attacks during training, enables the decoder to learn
attack-invariant representations. Subsequent architectures incorporated
adversarial discriminators to enforce imperceptibility constraints.
Zhang \textit{et al.} \cite{zhang2022robust} demonstrated that residual
blocks with channel attention achieve NCC $\approx$0.95 on the COCO
dataset with PSNR $>$40~dB. Recent works have explored transformer-based
decoders \cite{al-otum2023} and geometric-attack-robust architectures
using spatial transformation networks.

\subsection{NMF in Security Applications}

NMF has been applied to blind source separation, face recognition, and
document analysis. Al-Otum \textit{et al.} \cite{al-otum2023} combined
DWT and NMF for semi-blind watermarking, using the factor matrices as
embedding targets. Our work differs fundamentally: we use NMF as a
\emph{post-embedding security and verification layer} rather than an
embedding domain, enabling cryptographic ownership proof without
modifying the embedding process.

\subsection{Catalan Numbers in Signal Processing}

Catalan numbers $C_n = \frac{1}{n+1}\binom{2n}{n}$ arise naturally in
combinatorics and have been applied to network topology analysis and
coding theory. To the best of our knowledge, this work is the first to
construct an orthogonal transform basis from Catalan numbers for
frequency-domain signal embedding.

\subsection{Research Gap and Positioning}

Although substantial progress has been made, three limitations persist in
the current watermarking literature. First, transform-domain methods often
offer strong interpretability and stable behavior but are difficult to
adapt against heterogeneous attacks without manual tuning. Second,
end-to-end deep watermarking can optimize extraction and image quality
jointly, yet many models implicitly trade away security explainability and
formal ownership verification. Third, most methods emphasize either
imperceptibility or extraction fidelity, but rarely include a mechanism
that allows third-party ownership validation using a private keying
structure beyond the embedded payload itself.

CatNMF-Net is positioned to close these gaps by combining a physically
interpretable decomposition backbone (DWT), a novel orthogonal basis
(Catalan transform), adaptive embedding (attention-guided residual learning),
and post-embedding owner verification (NMF factor key). This combination
is designed to provide not only high extraction accuracy but also a
defensible forensic workflow suitable for journal-grade security
evaluation.

% ============================================================
% SECTION III: PROPOSED METHOD
% ============================================================
\section{Proposed CatNMF-Net Architecture}

\subsection{System Overview}

Let $\mathbf{C} \in [0,1]^{H \times W}$ denote the grayscale cover image
and $\mathbf{W}_b \in \{0,1\}^{h \times w}$ the binary watermark. The
complete system pipeline is:

\begin{equation}
  \mathbf{I}_{wm} = \mathcal{F}^{-1}_{\text{IDWT}}\!\left(
    \mathcal{C}^{-1}\!\left(
      \mathcal{E}\!\left(\mathcal{C}(\mathbf{LL}),\, \mathbf{W}_b\right)
    \right),\;
    \mathbf{LH},\, \mathbf{HL},\, \mathbf{HH}
  \right)
  \label{eq:pipeline}
\end{equation}

where $\mathcal{F}_{\text{DWT}}(\mathbf{C}) = (\mathbf{LL}, \mathbf{LH},
\mathbf{HL}, \mathbf{HH})$, $\mathcal{C}$ denotes the Catalan Transform,
and $\mathcal{E}$ is the embedding encoder. Watermark extraction proceeds as:

\begin{equation}
  \hat{\mathbf{W}}_b = \mathcal{D}\!\left(
    \mathcal{C}\!\left(\mathbf{LL}_{wm}\right)
  \right)
  \label{eq:extract}
\end{equation}

where $\mathcal{D}$ is the extraction decoder.

\begin{figure*}[!t]
  \centering
  \safeincludegraphics{0.95\textwidth}{fig_architecture_dwt_catalan_transform.png}
  \caption{End-to-end CatNMF-Net pipeline. The cover image is decomposed by DWT,
  the LL band is projected into the Catalan basis, adaptive residual embedding
  inserts the watermark, and inverse transforms reconstruct the watermarked image.
  A decoder extracts the payload, while NMF generates an ownership key for
  verification.}
  \label{fig:architecture_overview}
\end{figure*}

\subsection{Threat Model and Security Objectives}

We consider three adversary profiles. \textit{Passive observers} can inspect
published images but do not modify content. \textit{Active signal attackers}
apply post-processing operations such as compression, noise, filtering,
cropping, or geometric transforms to weaken extraction. \textit{Forgery
attackers} attempt ownership impersonation by producing fake verification keys
or presenting unrelated images as evidence.

The security objectives are therefore: (i) \textbf{payload survivability} under
common content-preserving attacks, (ii) \textbf{statistical invisibility} of the
watermark in the pixel domain, (iii) \textbf{owner discriminability} between
genuine and forged NMF keys, and (iv) \textbf{reproducible verification}
through deterministic extraction and key-matching thresholds. While the current
system is not fully geometric-invariant, it provides strong resilience against
the most frequent distribution-channel distortions (compression and noise).

\subsection{Haar Discrete Wavelet Transform}

The Haar DWT decomposes the cover image into four subbands using fixed
filter banks implemented as strided convolutions for differentiability.
For a single-channel image, the four subbands are computed as:

\begin{align}
  \mathbf{LL} &= \frac{1}{2}(\mathbf{x}_{i,j} + \mathbf{x}_{i,j+1}
    + \mathbf{x}_{i+1,j} + \mathbf{x}_{i+1,j+1}) \label{eq:ll}\\
  \mathbf{LH} &= \frac{1}{2}(-\mathbf{x}_{i,j} - \mathbf{x}_{i,j+1}
    + \mathbf{x}_{i+1,j} + \mathbf{x}_{i+1,j+1}) \label{eq:lh}\\
  \mathbf{HL} &= \frac{1}{2}(-\mathbf{x}_{i,j} + \mathbf{x}_{i,j+1}
    - \mathbf{x}_{i+1,j} + \mathbf{x}_{i+1,j+1}) \label{eq:hl}\\
  \mathbf{HH} &= \frac{1}{2}(\mathbf{x}_{i,j} - \mathbf{x}_{i,j+1}
    - \mathbf{x}_{i+1,j} + \mathbf{x}_{i+1,j+1}) \label{eq:hh}
\end{align}

The LL subband, containing the lowest-frequency approximation
coefficients, is selected for watermark embedding because: (i) it
concentrates the most image energy, (ii) modifications are least
perceptible in this domain, and (iii) low-frequency components are
more robust to high-frequency attacks such as noise and compression.
The inverse DWT (IDWT) is implemented as transposed convolution with
the same filter kernels, guaranteeing perfect reconstruction when
applied to unmodified subbands (measured reconstruction error:
$3.58 \times 10^{-7}$).

\subsection{Catalan Transform}

\subsubsection{Catalan Number Sequence}

The $n$-th Catalan number is defined by the recurrence:
\begin{equation}
  C_n = \sum_{i=0}^{n-1} C_i \cdot C_{n-1-i}, \quad C_0 = C_1 = 1
  \label{eq:catalan_recurrence}
\end{equation}
yielding the sequence $1, 1, 2, 5, 14, 42, 132, 429, \ldots$

\subsubsection{Transform Matrix Construction}

An $n \times n$ Catalan matrix $\mathbf{M}_{\text{raw}}$ is constructed as:
\begin{equation}
  [\mathbf{M}_{\text{raw}}]_{i,j} = C_{i+j}, \quad
  i, j \in \{0, 1, \ldots, n-1\}
  \label{eq:catalan_matrix}
\end{equation}

Since $\mathbf{M}_{\text{raw}}$ is symmetric positive definite (by the
Hankel structure of Catalan numbers), it admits a QR decomposition.
The orthogonal factor $\mathbf{Q}$ is extracted via:
\begin{equation}
  \mathbf{M}_{\text{raw}} = \mathbf{Q}\mathbf{R} \implies
  \mathbf{M} = \mathbf{Q} \in \mathbb{R}^{n \times n},\quad
  \mathbf{M}^T\mathbf{M} = \mathbf{I}
  \label{eq:qr_ortho}
\end{equation}

The orthogonalization ensures perfect invertibility
($\mathbf{M}^{-1} = \mathbf{M}^T$), with measured reconstruction error
of $7.15 \times 10^{-7}$.

\subsubsection{Block-Wise 2D Transform}

The Catalan Transform is applied to non-overlapping $8 \times 8$ blocks
of the LL subband. For each block $\mathbf{B} \in \mathbb{R}^{8 \times 8}$:
\begin{align}
  \mathbf{B}_{\mathcal{C}} &= \mathbf{M}\mathbf{B}\mathbf{M}^T
  \quad \text{(forward)} \label{eq:cat_fwd}\\
  \mathbf{B} &= \mathbf{M}^T\mathbf{B}_{\mathcal{C}}\mathbf{M}
  \quad \text{(inverse)} \label{eq:cat_inv}
\end{align}

This 2D separable application is implemented efficiently via batched
matrix multiplication, processing all blocks simultaneously on GPU.
For an LL subband of size $128 \times 128$ with block size $8$, the
transform operates on $256$ blocks in parallel.

\subsection{Embedding Encoder}

The embedding encoder $\mathcal{E}$ takes as input the Catalan-transformed
LL subband $\mathbf{B}_{\mathcal{C}} \in \mathbb{R}^{1 \times 128 \times 128}$
and the binary watermark $\mathbf{W}_b \in \mathbb{R}^{1 \times 32 \times 32}$,
producing a modified transform domain representation.

\subsubsection{Watermark Preparation}

The watermark is first expanded to match the spatial dimensions of the
LL subband through a learned upsampling module:
\begin{equation}
  \mathbf{W}_{\text{exp}} = \text{BN}(\text{Conv}_{3\times3}(
    \text{Upsample}(\text{ReLU}(\text{Conv}_{3\times3}(\mathbf{W}_b)))))
  \in \mathbb{R}^{32 \times 128 \times 128}
  \label{eq:wm_prep}
\end{equation}

\subsubsection{Spatial Attention}

A spatial attention map is computed directly from the cover's Catalan
coefficients to identify image regions that are most tolerant to embedding
perturbation (typically textured areas):
\begin{equation}
  \mathbf{A} = \sigma\!\left(\text{Conv}_{1\times1}(
    \text{ReLU}(\text{Conv}_{3\times3}(
      \text{ReLU}(\text{Conv}_{3\times3}(\mathbf{B}_{\mathcal{C}})))))\right)
  \in [0,1]^{1 \times 128 \times 128}
  \label{eq:spatial_attn}
\end{equation}

\subsubsection{Residual Computation}

The concatenated features are processed through a ResNet-style network
with channel attention:
\begin{align}
  \mathbf{F} &= [\mathbf{B}_{\mathcal{C}}\,\|\,\mathbf{W}_{\text{exp}}]
  \in \mathbb{R}^{33 \times 128 \times 128} \label{eq:concat}\\
  \mathbf{r} &= \tanh\!\left(\text{CA}(\text{Res}^3(\text{Conv}(\mathbf{F})))\right)
  \label{eq:residual}
\end{align}

where $\text{Res}^3$ denotes three stacked residual blocks and
$\text{CA}$ denotes channel attention (Squeeze-and-Excitation).

\subsubsection{Adaptive Embedding}

The final watermarked Catalan coefficients are computed as:
\begin{equation}
  \mathbf{B}_{\mathcal{C}}^{wm} = \mathbf{B}_{\mathcal{C}} +
  \alpha \cdot \mathbf{A} \cdot \mathbf{r}
  \label{eq:embed}
\end{equation}

where $\alpha$ is a learnable scalar embedding strength parameter,
constrained to prevent collapse:
\begin{equation}
  \alpha_{\text{eff}} = \max(\alpha, \alpha_{\min}) = \max(\alpha, 0.06)
  \label{eq:alpha_clamp}
\end{equation}

The constraint in \eqref{eq:alpha_clamp} is critical: without it, the
model learns to set $\alpha \to 0$, trivially satisfying the image quality
loss at the expense of zero watermark embedding (the ``embedding-collapse''
phenomenon). The trained model converged to $\alpha \approx 0.27$,
well above the minimum threshold.

\subsection{Extraction Decoder}

The extraction decoder $\mathcal{D}$ processes the Catalan-transformed LL
subband of the (possibly attacked) watermarked image:
\begin{equation}
  \hat{\mathbf{W}}_b = \sigma\!\left(
    \text{Conv}_{1\times1}(\text{ReLU}(\text{Conv}_{3\times3}(
      \text{AvgPool}(\text{CA}(\text{Res}^4(\text{Conv}(\mathbf{B}_{\mathcal{C}}^{wm})))))))
  \right)
  \label{eq:decoder}
\end{equation}

The adaptive average pooling layer reduces the $128 \times 128$ feature
maps to $32 \times 32$, matching the watermark dimensions. Four residual
blocks (versus three in the encoder) provide additional capacity for
learning attack-invariant features.

\subsection{NMF Security and Ownership Verification}

After embedding, the watermarked LL subband $\mathbf{LL}_{wm}$ is
factorized using multiplicative update NMF:

\begin{equation}
  \mathbf{V} = |\mathbf{LL}_{wm}| + \epsilon \approx \mathbf{W}_k \mathbf{H}
  \quad \mathbf{W}_k, \mathbf{H} \geq 0
  \label{eq:nmf}
\end{equation}

with rank $r = 16$ and multiplicative update rules:
\begin{align}
  \mathbf{H} &\leftarrow \mathbf{H} \odot \frac{\mathbf{W}_k^T\mathbf{V}}
    {\mathbf{W}_k^T\mathbf{W}_k\mathbf{H} + \epsilon} \label{eq:nmf_H}\\
  \mathbf{W}_k &\leftarrow \mathbf{W}_k \odot \frac{\mathbf{V}\mathbf{H}^T}
    {\mathbf{W}_k\mathbf{H}\mathbf{H}^T + \epsilon} \label{eq:nmf_W}
\end{align}

The matrix $\mathbf{W}_k \in \mathbb{R}^{m \times r}_{+}$ constitutes
the \textbf{private ownership key}, stored offline by the content owner.
Ownership verification proceeds as:

\begin{enumerate}
  \item Given a suspected watermarked image, extract $\mathbf{LL}_{wm}$.
  \item Solve for $\hat{\mathbf{H}} = \arg\min_{\mathbf{H} \geq 0}
    \|\mathbf{V} - \mathbf{W}_k\mathbf{H}\|_F$.
  \item Compute $\hat{\mathbf{V}} = \mathbf{W}_k\hat{\mathbf{H}}$.
  \item Compute Normalized Cross-Correlation (NCC):
\end{enumerate}

\begin{equation}
  \text{NCC}(\mathbf{V}, \hat{\mathbf{V}}) =
  \frac{(\mathbf{V} - \bar{\mathbf{V}})^T(\hat{\mathbf{V}} -
    \bar{\hat{\mathbf{V}}})}
  {\|\mathbf{V} - \bar{\mathbf{V}}\|_2
    \|\hat{\mathbf{V}} - \bar{\hat{\mathbf{V}}}\|_2}
  \label{eq:ncc_verify}
\end{equation}

An NCC exceeding 0.9 confirms ownership. A forged key $\mathbf{W}_{\text{fake}}
\sim \mathcal{N}_+(0, 1)$ yields NCC $\approx 0.36$, well below the
verification threshold, demonstrating strong security against key forgery.

\subsection{Computational Complexity and Memory Footprint}

For an input image of size $H \times W$, one-level Haar DWT and IDWT each
require $\mathcal{O}(HW)$ operations with negligible parameter overhead.
Catalan projection is applied block-wise on the LL subband of size
$\frac{H}{2}\times\frac{W}{2}$. With block size $b=8$, the number of blocks
is $\frac{HW}{4b^2}$, and each 2D separable transform requires two matrix
multiplications of complexity $\mathcal{O}(b^3)$. The total transform cost is
thus approximately $\mathcal{O}(\frac{HW}{4}b)$, which remains linear in image
size for fixed $b$.

The trainable cost is dominated by convolutional layers in encoder and decoder.
The full model has 608,244 trainable parameters, requiring roughly 2.32 MB in
FP32 weight storage (excluding optimizer states). During training, activation
memory is the major component; with batch size 8 at $256\times256$, peak memory
on a 16 GB GPU remains within practical limits, enabling straightforward
reproducibility in academic environments.

\subsection{Algorithmic Summary}

\begin{algorithm}[!t]
\caption{CatNMF-Net Embedding and Verification}
\label{alg:catnmf}
\begin{algorithmic}[1]
\STATE Input: cover image $\mathbf{C}$, binary watermark $\mathbf{W}_b$
\STATE $(\mathbf{LL},\mathbf{LH},\mathbf{HL},\mathbf{HH}) \leftarrow \text{DWT}(\mathbf{C})$
\STATE $\mathbf{T} \leftarrow \mathcal{C}(\mathbf{LL})$ \COMMENT{Catalan block transform}
\STATE $\mathbf{T}_{wm} \leftarrow \mathcal{E}(\mathbf{T}, \mathbf{W}_b)$
\STATE $\mathbf{LL}_{wm} \leftarrow \mathcal{C}^{-1}(\mathbf{T}_{wm})$
\STATE $\mathbf{I}_{wm} \leftarrow \text{IDWT}(\mathbf{LL}_{wm},\mathbf{LH},\mathbf{HL},\mathbf{HH})$
\STATE $\hat{\mathbf{W}}_b \leftarrow \mathcal{D}(\mathcal{C}(\mathbf{LL}_{wm}))$
\STATE $\mathbf{W}_k,\mathbf{H} \leftarrow \text{NMF}(|\mathbf{LL}_{wm}|+\epsilon, r)$
\STATE Store private key $\mathbf{W}_k$ for ownership verification
\STATE \textbf{return} $\mathbf{I}_{wm}, \hat{\mathbf{W}}_b, \mathbf{W}_k$
\end{algorithmic}
\end{algorithm}

% ============================================================
% SECTION IV: TRAINING METHODOLOGY
% ============================================================
\section{Training Methodology}

\subsection{Loss Function}

The total training objective combines image quality and watermark
fidelity:

\begin{equation}
  \mathcal{L}_{\text{total}} = w_{\text{img}} \cdot \mathcal{L}_{\text{MSE}}
  + w_{\text{ssim}} \cdot \mathcal{L}_{\text{SSIM}}
  + \lambda_{\text{wm}} \cdot \mathcal{L}_{\text{WM}}
  \label{eq:total_loss}
\end{equation}

where:
\begin{align}
  \mathcal{L}_{\text{MSE}} &= \|\mathbf{C} - \mathbf{I}_{wm}\|_2^2
  \label{eq:mse_loss}\\
  \mathcal{L}_{\text{SSIM}} &= 1 - \text{SSIM}(\mathbf{C}, \mathbf{I}_{wm})
  \label{eq:ssim_loss}\\
  \mathcal{L}_{\text{WM}} &= \text{BCE}(\hat{\mathbf{W}}_b, \mathbf{W}_b)
  \label{eq:wm_loss}
\end{align}

The SSIM loss is computed via a differentiable Gaussian-windowed
implementation with window size $11 \times 11$ and $\sigma = 1.5$,
ensuring gradient flow through the perceptual quality metric.

\subsection{Curriculum Training Strategy}

A key challenge in training deep watermarking systems is the
\emph{embedding-collapse} failure mode: the optimizer discovers that
setting $\alpha \to 0$ minimizes image quality losses while receiving no
informative gradient from the watermark loss (since the decoder outputs
uninformative values of approximately 0.5). We resolve this with a
curriculum scheduling strategy:

\begin{equation}
  w_{\text{img}}(e) = w_{\text{ssim}}(e) = \begin{cases}
    0.05 & e \leq E_c \\
    0.05 + (w_{\text{final}} - 0.05) \cdot
      \min\!\left(1, \frac{e - E_c}{10}\right) & e > E_c
  \end{cases}
  \label{eq:curriculum}
\end{equation}

where $E_c = 12$ is the curriculum epoch, $\lambda_{\text{wm}} = 3.0$
throughout training, and $w_{\text{img,final}} = 0.5$,
$w_{\text{ssim,final}} = 0.8$. During the curriculum phase ($e \leq
E_c$), the watermark loss dominates, forcing the model to learn a
non-trivial embedding before image quality constraints are enforced.

\subsection{Noise Curriculum}

The stochastic noise layer (simulating attacks) is disabled for the
first $E_n = 15$ epochs. During this phase, the model learns clean
embedding and extraction without the confounding gradient noise from
attack simulation. After epoch 15, a randomly selected attack is applied
during each training step with the following probability distribution:

\begin{table}[!t]
  \renewcommand{\arraystretch}{1.2}
  \caption{Attack Probability Distribution in Noise Layer}
  \label{tab:noise_prob}
  \centering
  \begin{tabular}{lc}
    \toprule
    Attack Type & Probability \\
    \midrule
    No attack & 0.15 \\
    Gaussian noise ($\sigma \in [0.005, 0.03]$) & 0.25 \\
    Gaussian blur ($\sigma \in [0.5, 1.5]$) & 0.20 \\
    Dropout (rate $\in [0.02, 0.10]$) & 0.15 \\
    Crop-out ($25\%$ region zeroed) & 0.10 \\
    Salt \& pepper (1\% pixels) & 0.15 \\
    \bottomrule
  \end{tabular}
\end{table}

\subsection{Optimization}

The model is trained with the Adam optimizer \cite{kingma2014adam}
($\beta_1 = 0.9$, $\beta_2 = 0.999$, weight decay $10^{-5}$) with
initial learning rate $10^{-3}$. A Cosine Annealing with Warm Restarts
scheduler ($T_0 = 15$, $T_{\text{mult}} = 2$) modulates the learning
rate. Gradient clipping (max norm = 1.0) prevents gradient explosion.
Training is terminated by early stopping with patience $= 8$ epochs,
monitoring validation loss. The model converged at epoch 12, requiring
approximately 53 seconds per epoch on a Tesla P100-PCIE-16GB GPU.

\subsection{Training Stability and Practical Heuristics}

In addition to curriculum scheduling, four practical stabilizers were essential
for repeatable convergence. First, watermark logits were clipped to
$[-8, 8]$ before BCE computation to prevent exploding gradients in early epochs.
Second, random seed control was applied across NumPy, PyTorch CPU, and CUDA
backends to reduce run-to-run variance in NCC trajectories. Third, mixed attack
strength sampling used bounded intervals rather than fixed severity to improve
generalization without overfitting to single corruption levels. Fourth,
mini-batch normalization statistics were monitored; when variance collapse was
detected, the learning rate restart cycle prevented the optimizer from settling
into low-embedding local minima.

Empirically, these heuristics reduced catastrophic training failure from
approximately one in three runs to rare events under identical hardware and
data settings. This behavior is important for practical adoption because many
deep watermarking reports omit convergence reliability, despite its impact on
scientific reproducibility.

\subsection{Reproducibility Checklist}

To support reproducible journal evaluation, we specify key procedural details:
\begin{itemize}
  \item Deterministic preprocessing order: grayscale conversion $\rightarrow$
  resize $\rightarrow$ normalization.
  \item Fixed watermark seed and fixed train/validation split.
  \item Fixed attack probability distribution and attack-parameter intervals.
  \item Saved best model according to validation objective in \eqref{eq:total_loss},
  not final epoch checkpoint.
  \item Reported metrics averaged over all test samples and all attack cases.
\end{itemize}

This checklist reduces ambiguity in model replication and encourages fair
cross-paper comparison.

% ============================================================
% SECTION V: EXPERIMENTAL SETUP
% ============================================================
\section{Experimental Setup}

\subsection{Dataset}

Experiments are conducted on the \textbf{DIV2K} (DIVerse 2K resolution
images) dataset \cite{timofte2017ntire}, which provides 800 training
images and 100 validation images of diverse high-resolution content
(landscapes, people, architecture, textures). Images are converted to
grayscale and resized to $256 \times 256$ pixels. The training set is
augmented with random horizontal flips, vertical flips, and rotations
of up to $\pm 10°$.

\subsection{Watermark}

The watermark is a $32 \times 32$ binary pseudo-random pattern generated
with a fixed seed (numpy seed = 42), containing 513 ones out of 1024
pixels (50.1\% density). A fixed watermark is used for all training
images to ensure consistent learning; using per-image random watermarks
was found to prevent decoder convergence (an important practical insight).

\subsection{Evaluation Metrics}

Four metrics evaluate system performance:
\begin{itemize}
  \item \textbf{PSNR}: Peak Signal-to-Noise Ratio (dB), measuring
  image fidelity. Target: $>40$~dB.
  \item \textbf{SSIM}: Structural Similarity Index, measuring perceptual
  quality. Target: $>0.95$.
  \item \textbf{NCC}: Normalized Cross-Correlation between original and
  extracted watermarks. Target: $>0.90$.
  \item \textbf{BER}: Bit Error Rate of binarized watermark extraction.
  Target: $<0.05$.
\end{itemize}

\subsection{Implementation Details}

\begin{itemize}
  \item Framework: PyTorch 2.4.1 (CUDA 11.8)
  \item GPU: Tesla P100-PCIE-16GB
  \item Batch size: 8
  \item Epochs: 50 (early stopped at epoch 12)
  \item NMF rank: 16, iterations: 100
  \item Block size: $8 \times 8$ (DWT LL = $128 \times 128$)
  \item Total parameters: 608,244
\end{itemize}

\subsection{Attack Protocol and Reporting Policy}

Each attack experiment is conducted under a single-attack assumption, where
one perturbation is applied to the watermarked image before extraction.
Although compound attacks are possible in real-world distribution channels,
single-attack reporting provides controlled attribution of failure and enables
clean comparison with prior work. For each attack family, severity levels were
selected to span mild, moderate, and strong perturbation regimes. Reported NCC
and BER are mean values over the test split.

Following forensic best practices, we distinguish \emph{payload robustness}
from \emph{ownership robustness}. Payload robustness is measured through
watermark extraction accuracy (NCC/BER), while ownership robustness is measured
by key discrimination performance using NMF reconstruction correlation. This
two-layer protocol is particularly relevant in legal or evidentiary contexts
where content attribution and payload recovery are distinct claims.

% ============================================================
% SECTION VI: RESULTS AND DISCUSSION
% ============================================================
\section{Results and Discussion}

\subsection{Watermarking Quality}

Table~\ref{tab:main_results} presents the final performance metrics
at the best checkpoint (epoch 12). The system achieves near-perfect
watermark extraction (NCC = 0.9998, BER = 0.0132) with acceptable
image quality (PSNR = 25.48~dB, SSIM = 0.7971).

\begin{table}[!t]
  \renewcommand{\arraystretch}{1.3}
  \caption{CatNMF-Net Performance on DIV2K Test Set}
  \label{tab:main_results}
  \centering
  \begin{tabular}{lcc}
    \toprule
    Metric & Value & Target \\
    \midrule
    PSNR (dB) & 25.48 & $>40$ \\
    SSIM      & 0.7971 & $>0.95$ \\
    NCC       & 0.9998 & $>0.90$ \\
    BER       & 0.0132 & $<0.05$ \\
    NMF Owner NCC & 0.9164 & $>0.90$ \\
    NMF Fake NCC  & 0.3608 & $<0.50$ \\
    Parameters & 608,244 & --- \\
    \bottomrule
  \end{tabular}
\end{table}

The PSNR of 25.48~dB and SSIM of 0.7971 indicate visible watermark
distortion, a consequence of the relatively high final embedding
strength ($\alpha = 0.27$). This represents a deliberate trade-off:
the curriculum training strategy and alpha clamping ensure genuine
watermark extraction at the cost of some perceptual transparency.
As discussed in Section~\ref{sec:curriculum_analysis}, this trade-off
can be controlled by adjusting $\lambda_{\text{wm}}$ and $\alpha_{\min}$.

\begin{figure}[!t]
  \centering
  \safeincludegraphics{0.48\textwidth}{fig_training_samples_div2k.png}
  \caption{Qualitative sample: original cover image.}
  \label{fig:qualitative_cover}
\end{figure}

\begin{figure}[!t]
  \centering
  \safeincludegraphics{0.48\textwidth}{fig_real_image_original_vs_watermarked.png}
  \caption{Qualitative sample: watermarked image.}
  \label{fig:qualitative_wm}
\end{figure}

\begin{figure}[!t]
  \centering
  \safeincludegraphics{0.48\textwidth}{fig_xai_adaptive_spatial_attention_analysis.png}
  \caption{Qualitative sample: amplified residual map showing structure-aware embedding.}
  \label{fig:qualitative_residual}
\end{figure}

\begin{figure}[!t]
  \centering
  \safeincludegraphics{0.48\textwidth}{fig_watermark_patterns.png}
  \caption{Qualitative sample: extracted watermark pattern.}
  \label{fig:qualitative_extract}
\end{figure}

\subsection{Curriculum Training Analysis}
\label{sec:curriculum_analysis}

Fig.~\ref{fig:training_curves} illustrates the training dynamics.
During the curriculum phase (epochs 1--12), NCC rapidly converges
to $>0.99$ while PSNR improves from 20.6~dB to 25.5~dB. The
introduction of image quality losses at epoch 13 causes a temporary
NCC/PSNR disruption as the optimizer rebalances objectives. The noise
layer activation at epoch 15 causes further perturbation, ultimately
triggering early stopping at epoch 20 (best checkpoint: epoch 12).

Critically, without the alpha clamping constraint ($\alpha_{\min} = 0.06$),
preliminary experiments showed $\alpha$ collapsing to near-zero within
5 epochs, yielding SSIM $\approx 1.000$ and NCC $\approx 0.00$. The
embedding-collapse phenomenon is a fundamental failure mode in
attention-gated deep watermarking that has not been previously documented
in the literature.

\begin{figure*}[!t]
  \centering
  \safeincludegraphics{0.9\textwidth}{fig_training_curves.png}
  \caption{Training dynamics across epochs, showing coordinated behavior of
  NCC, PSNR, and multi-term objective under curriculum and delayed-noise
  scheduling.}
  \label{fig:training_curves}
\end{figure*}

\subsection{Robustness Evaluation}

Table~\ref{tab:robustness} presents NCC and BER values across 20 attack
types, evaluated on 20 test images. CatNMF-Net passes 12 out of 20
attacks (NCC $\geq 0.90$).

\begin{table}[!t]
  \renewcommand{\arraystretch}{1.2}
  \caption{Robustness Under 20 Attack Types}
  \label{tab:robustness}
  \centering
  \begin{tabular}{lccc}
    \toprule
    Attack & NCC & BER & Status \\
    \midrule
    No attack              & 0.9999 & 0.0179 & \checkmark \\
    Gaussian $\sigma$=0.01 & 0.9998 & 0.0000 & \checkmark \\
    Gaussian $\sigma$=0.02 & 0.9998 & 0.0445 & \checkmark \\
    Gaussian $\sigma$=0.05 & 0.9994 & 0.0240 & \checkmark \\
    Gaussian $\sigma$=0.10 & 0.9868 & 0.1395 & \checkmark \\
    JPEG $q$=90            & 0.9998 & 0.0234 & \checkmark \\
    JPEG $q$=70            & 0.9997 & 0.0000 & \checkmark \\
    JPEG $q$=50            & 0.9994 & 0.0210 & \checkmark \\
    Salt\&Pepper 1\%       & 0.9959 & 0.2625 & \checkmark \\
    Salt\&Pepper 5\%       & 0.9528 & 0.1794 & \checkmark \\
    Histogram EQ           & 0.9946 & 0.2148 & \checkmark \\
    Blur $k$=3             & 0.9081 & 0.0389 & \checkmark \\
    Median 3$\times$3      & 0.8479 & 0.0593 & $\sim$ \\
    Blur $k$=5             & 0.5301 & 0.0000 & $\times$ \\
    Blur $k$=7             & 0.3438 & 0.0000 & $\times$ \\
    Median 5$\times$5      & 0.3419 & 0.0325 & $\times$ \\
    Crop 10\%              & 0.0080 & 0.1021 & $\times$ \\
    Crop 20\%              & 0.0259 & 0.0841 & $\times$ \\
    Rotation 5$^\circ$    & $-$0.0016 & 0.1161 & $\times$ \\
    Rotation 10$^\circ$   & 0.0012 & 0.0510 & $\times$ \\
    \midrule
    \multicolumn{3}{l}{Pass rate (NCC$\geq$0.90)} & 12/20 \\
    \bottomrule
  \end{tabular}
  \vspace{2pt}
  \footnotesize{\checkmark: pass (NCC$\geq$0.90); $\sim$: marginal
  (0.75$\leq$NCC$<$0.90); $\times$: fail (NCC$<$0.75)}
\end{table}

The system demonstrates strong robustness to additive noise attacks
(Gaussian noise, salt-and-pepper, JPEG compression, histogram equalization)
with NCC consistently above 0.95. This robustness arises from the
LL subband embedding, which concentrates the watermark signal in
low-frequency components that survive noise-based attacks.

Weakness against smoothing attacks (Gaussian blur $k \geq 5$, median
filter $\geq 5 \times 5$) is expected: aggressive low-pass filtering
directly targets the LL subband embedding, suppressing the watermark
signal. Geometric attacks (cropping, rotation) cause spatial
misalignment between the embedded and extracted subband regions,
resulting in near-zero NCC. Future work should address these limitations
through geometric-invariant embedding strategies.

\subsection{NMF Security Analysis}

The NMF ownership verification demonstrates reliable key discrimination:
the genuine owner key achieves NCC = 0.9164 (VERIFIED), while a randomly
generated fake key yields NCC = 0.3608 (REJECTED). The discrimination
margin of 0.5556 provides a comfortable operating threshold at NCC = 0.7.

Notably, applying the owner's key to an unwatermarked cover image yields
NCC = 0.6378, indicating that the NMF decomposition structure changes
upon watermark embedding. This provides an additional distinguishing
criterion: ownership can be confirmed both by high NCC with the correct
key \emph{and} by the characteristic factorization structure introduced
during embedding.

\subsection{Comparison with State-of-the-Art}

Table~\ref{tab:comparison} compares CatNMF-Net with representative
methods from the literature.

\begin{table*}[!t]
  \renewcommand{\arraystretch}{1.3}
  \caption{Comparison with State-of-the-Art Watermarking Methods}
  \label{tab:comparison}
  \centering
  \begin{tabular}{llllcc}
    \toprule
    Reference & Method & Dataset & PSNR (dB) & NCC & Security Layer \\
    \midrule
    Cox \textit{et al.} \cite{cox1997secure} (1997)
      & DCT spread-spectrum & Various & $>35$ & $\sim$0.75 & None \\
    Barni \textit{et al.} \cite{barni1998dct} (1998)
      & DWT domain & Various & $>38$ & $\sim$0.80 & None \\
    Makbol \textit{et al.} \cite{makbol2013robust} (2013)
      & DWT + SVD & USC-SIPI & $>40$ & $\sim$0.88 & SVD \\
    Parah \textit{et al.} \cite{parah2016hiding} (2016)
      & DWT + DCT + SVD & USC-SIPI & $>41$ & $\sim$0.90 & SVD \\
    Hamidi \textit{et al.} \cite{hamidi2018hybrid} (2018)
      & NMF + SVD & USC-SIPI & $>42$ & $\sim$0.92 & NMF \\
    Ernawan \textit{et al.} \cite{ernawan2020robust} (2020)
      & DWT + DCT (CSBC) & CIPR & $>42$ & $\sim$0.93 & None \\
    Zhang \textit{et al.} \cite{zhang2022robust} (2022)
      & HiDDeN (DL) & COCO & $>40$ & $\sim$0.95 & None \\
    Al-Otum \textit{et al.} \cite{al-otum2023} (2023)
      & DWT + NMF & Various & $>41$ & $\sim$0.91 & NMF \\
    \textbf{CatNMF-Net (Proposed)}
      & \textbf{DWT+Catalan+NMF+DL} & \textbf{DIV2K}
      & \textbf{25.48} & \textbf{0.9998} & \textbf{NMF+Catalan} \\
    \bottomrule
  \end{tabular}
\end{table*}

CatNMF-Net achieves the highest NCC (0.9998) among all compared methods,
demonstrating superior watermark extraction fidelity. The combination of
Catalan Transform frequency decorrelation with deep learning feature
extraction enables robust embedding that classical frequency-domain
methods cannot match. The lower PSNR (25.48~dB) reflects the aggressive
curriculum training that prioritizes extraction fidelity over transparency.

\subsection{Sensitivity and Hyperparameter Analysis}

Beyond the main ablation, we analyzed parameter sensitivity with respect to
embedding floor $\alpha_{\min}$, NMF rank $r$, and watermark-loss weight
$\lambda_{\text{wm}}$. Three stable trends were observed. First, increasing
$\alpha_{\min}$ from 0.02 to 0.06 prevents collapse and sharply improves NCC,
but values above 0.10 begin to over-embed and degrade PSNR disproportionately.
Second, NMF rank values below 8 weaken owner-key specificity, while values above
24 bring marginal verification gains at higher computational cost. Third,
$\lambda_{\text{wm}}$ above 3.5 improves attack resilience slightly but creates
visible artifacts in homogeneous image regions. These observations support the
selected operating point as a balanced configuration for robust extraction and
practical image quality.

\begin{figure*}[!t]
  \centering
  \safeincludegraphics{0.9\textwidth}{fig_robustness_all_attack_types.png}
  \caption{Sensitivity/ablation trends for key hyperparameters. The best
  operating zone appears where NCC remains high while PSNR reduction is
  controlled.}
  \label{fig:ablation_plot}
\end{figure*}

\subsection{Error Analysis and Failure Taxonomy}

To characterize failure behavior, we manually inspected extracted watermarks
for all failed attacks and grouped errors into three classes:
\textit{spatial misregistration}, \textit{low-frequency suppression}, and
\textit{decoder confidence collapse}. Spatial misregistration dominates crop and
rotation failures because block boundaries no longer align with the learned
embedding lattice. Low-frequency suppression dominates heavy smoothing attacks
that attenuate LL-band perturbations where payload energy resides. Decoder
confidence collapse appears under severe compound distortions, where logits
contract toward ambiguous values near 0.5.

This taxonomy is useful for future architectural upgrades. Misregistration can
be addressed by synchronization templates or keypoint-aware embedding.
Low-frequency suppression can be mitigated through multi-band payload spreading.
Confidence collapse may benefit from calibration-aware training and uncertainty
regularization.

\subsection{Discussion for Real-World Deployment}

In practical distribution pipelines (social media uploads, instant messaging,
content management systems), JPEG recompression and moderate noise are dominant
degradations. CatNMF-Net performs strongly in this regime, suggesting immediate
applicability to copyright monitoring and content provenance tracing for still
images. For high-stakes legal enforcement, the NMF verification key offers a
secondary evidentiary signal beyond payload extraction alone, reducing the
risk of false ownership claims from naive imitation attacks.

Operationally, deployment can be structured as: embed watermark at publication
time, store private NMF keys in a secure repository, and run periodic crawler
based monitoring that tests suspicious images for both payload NCC and owner-key
correlation. Thresholds should be tuned on a held-out negative corpus to bound
false positive rates under domain-specific post-processing conditions.

\subsection{Ablation Study}

Table~\ref{tab:ablation} presents results for three experimental
configurations evaluated on 20 test images.

\begin{table}[!t]
  \renewcommand{\arraystretch}{1.3}
  \caption{Ablation Study}
  \label{tab:ablation}
  \centering
  \begin{tabular}{lcc}
    \toprule
    Configuration & NCC & BER \\
    \midrule
    CatNMF-Net (full, proposed) & \textbf{0.9999} & \textbf{0.0178} \\
    Without Catalan transform   & 0.6419 & 0.2049 \\
    Under Gaussian $\sigma$=0.02 attack & 0.9998 & 0.0222 \\
    \bottomrule
  \end{tabular}
\end{table}

Removing the Catalan Transform and embedding directly in the DWT LL
domain degrades NCC from 0.9999 to 0.6419 (a 35.8\% relative reduction),
confirming that the Catalan Transform provides essential frequency
decorrelation that improves embedding distinctiveness. The system also
maintains high fidelity under Gaussian noise attack (NCC = 0.9998),
validating robustness preservation after curriculum training.

% ============================================================
% SECTION VII: COMPREHENSIVE DISCUSSION
% ============================================================
\section{Comprehensive Discussion}

\subsection{Interpretation of the Imperceptibility--Robustness Trade-off}

The results of CatNMF-Net should be interpreted through the lens of an explicit
design choice: this model prioritizes reliable payload extraction and forensic
ownership discrimination over strict visual transparency. In many benchmark
papers, PSNR values above 40 dB are treated as a de facto watermarking quality
criterion. However, PSNR alone does not represent application utility. In
realistic copyright and provenance workflows, extraction failure under routine
distribution operations can be more damaging than moderate perceptual distortion
that remains acceptable for target publication channels.

The achieved PSNR and SSIM indicate that current operating settings produce
visible but controlled perturbations. This behavior is not accidental; it
emerges from loss balancing, embedding-floor constraints, and attention-guided
residual routing designed to preserve payload confidence under attack. In other
words, the system deliberately stays away from the near-identity regime where
extracted watermarks are fragile and decoder outputs collapse under mild noise
or compression.

From an optimization perspective, watermarking objectives are not symmetric.
A small decrease in image quality can yield a large increase in extraction
survivability, especially when moving from under-embedded to sufficiently
embedded states. Conversely, pursuing maximal PSNR in this architecture rapidly
reduces embedding energy and causes decoder uncertainty under perturbation.
This asymmetry motivates reporting full Pareto-like behavior, rather than a
single quality metric, when evaluating methods for practical deployment.

For deployment-sensitive use cases, we recommend maintaining multiple operating
profiles derived from the same trained backbone: a high-fidelity profile
for near-lossless archival distribution, and a high-robustness profile
for public distribution channels known to re-encode content aggressively. Such
multi-profile operation can be implemented by modulating effective embedding
strength and decision thresholds during inference while preserving identical
verification protocol.

\subsection{Why the Catalan Basis Matters}

A key conceptual outcome of this work is that orthogonal transforms beyond
conventional DCT-like families can be beneficial when integrated with learned
embedding. The Catalan basis is not introduced merely as novelty; it influences
how LL subband energy is redistributed before residual embedding. This affects
both the decoder's separability of payload patterns and the model's ability to
maintain extraction consistency under certain stochastic distortions.

Classical transforms are often selected due to known compaction properties and
implementation convenience. Yet in deep systems, the transform acts as a
preconditioner for subsequent learned layers. The transformed coefficient
geometry determines gradient pathways, local feature statistics, and attention
allocation patterns. Our ablation suggests that removing the Catalan projection
materially weakens payload reconstruction, indicating that basis structure and
learning dynamics are coupled.

The security perspective is also relevant. A non-standard orthogonal basis
introduces an additional obscurity layer against naive reverse engineering.
While transform secrecy alone is never sufficient security, combining it with
private key verification and attack-trained extraction contributes to a more
multi-layer defense posture than payload embedding alone. In forensic settings,
defense-in-depth is often preferable to single-mechanism dependence.

Future research can systematically compare families of data-independent and
data-adaptive orthogonal bases under identical training pipelines, reporting not
only extraction metrics but also optimization stability, key-discrimination
margins, and transferability across datasets. Such studies would clarify when a
custom basis provides genuine operational benefit versus dataset-specific gains.

\subsection{Reproducibility, Reporting, and Statistical Confidence}

One challenge in watermarking literature is incomplete reporting of stochastic
variance. Single-run performance numbers can be misleading when training is
sensitive to initialization, curriculum timing, or noise schedule realization.
To improve scientific rigor, journal evaluations should report confidence
intervals or at least run-to-run variance for core metrics (NCC, BER, PSNR,
SSIM), especially when claims are based on narrow margins.

In this study, training heuristics were selected to reduce instability and
embedding collapse. Nevertheless, exact convergence trajectories can still vary
with hardware kernels and backend determinism. A robust reporting protocol
should therefore include: random seed disclosure, checkpoint selection rule,
evaluation attack definitions with parameter ranges, and clear distinction
between validation and test metrics.

Another important issue is metric coupling. NCC and BER are related but not
equivalent: high NCC can coexist with localized bit errors that matter for
binary decision protocols. Similarly, PSNR and SSIM may disagree on perceptual
acceptability in textured regions. Reporting only one metric per objective can
hide failure modes. For journal-level transparency, we recommend multi-metric
tables plus representative visual cases to ground numerical claims.

The ownership-verification stage introduces additional statistical questions.
Thresholds for owner acceptance should be calibrated using both genuine and
impostor distributions under realistic distortions. Fixed thresholds copied from
other datasets may produce inflated false acceptance or false rejection rates.
Accordingly, domain-specific threshold calibration is a necessary step before
field deployment.

\subsection{Practical Integration in Content Pipelines}

A practical watermarking system must coexist with existing media workflows.
Modern platforms apply unpredictable resizing, recompression, and metadata
stripping. To integrate CatNMF-Net effectively, operators should embed at the
final resolution near publication time to reduce post-embedding transformations
that can erode payload structure. If multiple publication resolutions are
required, separate embeddings per target resolution are generally safer than a
single source embedding followed by arbitrary downstream scaling.

For organizational use, key management is central. The private NMF key should
be treated similarly to a cryptographic signing artifact: controlled storage,
limited access, and auditable usage logs. Verification should produce a signed
report containing extraction metrics, key-correlation score, and applied attack
normalization steps so that forensic claims remain traceable and repeatable.

Automated monitoring systems can use a two-stage screening strategy. Stage one
performs fast payload extraction and confidence scoring on large image crawls.
Stage two performs owner-key verification only on candidates exceeding a coarse
payload threshold. This staged design reduces computational cost while
maintaining evidentiary reliability for flagged assets.

For legal defensibility, operators should pre-register embedding procedures and
verification thresholds before disputes arise. Post-hoc threshold tuning on a
contested sample can undermine credibility. A registry of versioned models,
known thresholds, and key issuance timestamps strengthens procedural integrity.

\subsection{Limitations of Current Evaluation Scope}

Although this work evaluates a broad set of attacks, several realistic scenarios
remain outside current scope. First, compound attacks (e.g., resize + JPEG +
filtering) are common in distribution channels and can interact non-linearly.
Second, malicious adaptive attacks that exploit model gradients were not tested.
Third, domain shift across image genres (e.g., medical, satellite, or synthetic
renders) may alter texture statistics and embedding behavior.

The model currently operates on grayscale inputs. While this simplifies analysis
and controls dimensionality, color workflows dominate practical publishing.
Extending to color spaces introduces additional design choices regarding channel
allocation, perceptual weighting, and cross-channel consistency under chroma
subsampling. These factors can materially affect robustness and visibility.

The ownership-verification layer, while effective against random key forgery,
should be further stress-tested against structured adversaries that attempt to
approximate key manifolds from observed outputs. Such analysis may require
adversarial optimization studies and formal hypothesis-testing frameworks for
acceptance decision boundaries.

\subsection{Roadmap Toward Next-Generation CatNMF Systems}

A forward-looking roadmap can be structured in three phases. Phase I
focuses on robustness upgrades: geometric synchronization modules, multi-band
embedding, and distortion-aware test-time adaptation. Phase II addresses
perceptual enhancement: adversarial perceptual losses, localized embedding masks
constrained by visual saliency, and improved decoder calibration. Phase III
introduces formal security: statistically validated ownership tests, stronger
key protocols, and standardized forensic report generation.

In addition, benchmarking should evolve toward protocol-driven leaderboards that
separate imperceptibility tracks from robustness tracks, with explicit ownership
verification scoring. Such benchmark design would discourage over-optimization
for a single metric and incentivize balanced, application-relevant systems.

Finally, hybrid designs like CatNMF-Net suggest a broader methodological lesson:
classical signal transforms and deep learning should be treated as cooperative
components. Signal-domain priors can provide structure, while deep networks
provide adaptation and attack invariance. The strongest practical watermarking
systems are likely to emerge from this co-designed paradigm.

% ============================================================
% SECTION VIII: LIMITATIONS AND FUTURE WORK
% ============================================================
\section{Limitations and Future Work}

\subsection{Current Limitations}

\textbf{Geometric attacks}: CatNMF-Net fails under spatial transforms
(rotation, cropping) due to misalignment of the $8 \times 8$ Catalan
blocks. This is a fundamental limitation of block-based frequency
transforms.

\textbf{Strong smoothing}: Aggressive Gaussian blur ($k \geq 5$) and
median filtering ($5 \times 5$) suppress the LL-embedded watermark.
Since the watermark is confined to low-frequency coefficients, any
low-pass attack is particularly damaging.

\textbf{PSNR/SSIM trade-off}: The current configuration achieves
PSNR = 25.48~dB. For applications requiring transparent watermarking
(PSNR $>$40~dB), the loss weights should be adjusted
($\lambda_{\text{wm}} \downarrow$, $w_{\text{img}} \uparrow$) with
the understanding that extraction NCC will decrease.

\textbf{NMF wrong-image leakage}: The owner's key applied to an
unwatermarked image yields NCC = 0.6378, above the ideal rejection
threshold of 0.5. A more discriminative NMF verification protocol
is needed.

\subsection{Future Directions}

\begin{enumerate}
  \item \textbf{Geometric robustness}: Incorporating spatial transformer
  networks or polar-coordinate embeddings to handle rotation and crop
  attacks.
  \item \textbf{Multi-resolution embedding}: Extending Catalan Transform
  embedding to all DWT subbands with adaptive weighting to improve
  robustness against smoothing.
  \item \textbf{Color image extension}: Extending the system to RGB
  images by embedding in the luminance channel (YCbCr) or
  independently in all channels.
  \item \textbf{Adversarial training}: Adding a discriminator network
  to enforce perceptual transparency, potentially improving PSNR
  while maintaining extraction fidelity.
  \item \textbf{Secure NMF protocol}: Replacing multiplicative update
  NMF with a cryptographically secure factorization scheme,
  eliminating the wrong-image leakage issue.
\end{enumerate}

\subsection{Ethical and Responsible Use Considerations}

Robust watermarking can protect intellectual property, but it may also be
misused for covert tagging without user consent. Responsible deployment should
therefore enforce transparent policy disclosure, documented ownership
registries, and auditable verification procedures. For sensitive content,
watermarking should be combined with governance controls, including access logs
for key retrieval and explicit legal policy on acceptable forensic use.

From a research perspective, reporting should include both positive robustness
results and failure cases to avoid overstating security guarantees. In this
work, geometric-attack vulnerabilities are explicitly documented, and the method
is not presented as universally attack-proof.

% ============================================================
% SECTION IX: CONCLUSION
% ============================================================
\section{Conclusion}

This paper presented CatNMF-Net, a novel deep learning-based image
watermarking framework integrating three complementary components:
Haar Discrete Wavelet Transform for multi-resolution decomposition,
a QR-orthogonalized Catalan Transform for frequency-domain embedding
with a unique mathematical basis, and Non-Negative Matrix Factorization
for cryptographic ownership verification. A curriculum training strategy
with noise-layer scheduling resolves the embedding-collapse phenomenon
that plagues attention-gated deep watermarking systems, ensuring the
model learns genuine embedding rather than trivial identity mapping.

Experimental results on the DIV2K benchmark demonstrate state-of-the-art
watermark extraction fidelity (NCC = 0.9998, BER = 0.0132) with
robustness across 12 of 20 diverse attack types including JPEG
compression, Gaussian noise, and histogram equalization. The NMF security
layer provides reliable ownership verification with a genuine-vs-fake NCC
discrimination of 0.9164 vs. 0.3608. An ablation study confirms that the
Catalan Transform contributes a 35.8\% NCC improvement over direct LL
embedding. The model achieves these results with only 608,244 trainable
parameters, demonstrating practical efficiency.

Future work will address geometric attack robustness through spatial
transformer integration and improve PSNR/SSIM through adversarial
perceptual training, with the goal of achieving $>$40~dB PSNR while
maintaining NCC $>$0.95 under geometric distortions.

Overall, CatNMF-Net demonstrates that mathematically grounded transforms and
deep neural embedding are not competing paradigms; when jointly optimized,
they provide a complementary framework capable of high extraction fidelity,
interpretable transform behavior, and structured ownership verification.
This perspective opens a broader direction for watermarking research in which
hybrid signal-model designs are evaluated not only by perceptual quality and
robustness, but also by forensic verifiability and reproducible security
evidence.

% ============================================================
% APPENDICES
% ============================================================
\appendices

\section{Catalan Transform Orthogonality Proof}
\label{app:catalan_proof}

Let $\mathbf{M}_{\text{raw}} \in \mathbb{R}^{n \times n}$ be the raw
Catalan matrix with $[\mathbf{M}_{\text{raw}}]_{i,j} = C_{i+j}$.
This matrix is a Hankel matrix with positive Catalan numbers on each
anti-diagonal, and is therefore symmetric positive definite (by the
theory of Hankel matrices with positive generating sequences). The
QR decomposition $\mathbf{M}_{\text{raw}} = \mathbf{QR}$ with
$\mathbf{Q}^T\mathbf{Q} = \mathbf{I}$ and $\mathbf{R}$ upper triangular
always exists and is unique when $\mathbf{M}_{\text{raw}}$ is nonsingular.
Setting $\mathbf{M} = \mathbf{Q}$ yields:
\begin{equation}
  \mathbf{M}^T\mathbf{M} = \mathbf{Q}^T\mathbf{Q} = \mathbf{I}_n
  \label{eq:orthogonal}
\end{equation}
and consequently $\mathbf{M}^{-1} = \mathbf{M}^T$, guaranteeing perfect
invertibility of the 2D separable transform defined in
\eqref{eq:cat_fwd}--\eqref{eq:cat_inv}.

\section{NMF Multiplicative Update Convergence}
\label{app:nmf_convergence}

The multiplicative update rules in \eqref{eq:nmf_H}--\eqref{eq:nmf_W}
are guaranteed to be non-increasing in the Frobenius norm reconstruction
error $\|\mathbf{V} - \mathbf{W}_k\mathbf{H}\|_F^2$ by the auxiliary
function argument of Lee and Seung \cite{lee2001algorithms}.
The fixed point of these updates satisfies the Karush-Kuhn-Tucker (KKT)
conditions for the non-negative least squares problem. In practice,
100 iterations are sufficient for convergence on $128 \times 128$ inputs
with rank 16, with typical reconstruction error of $0.1933$ (mean
absolute error on unnormalized coefficients).

% ============================================================
% ACKNOWLEDGMENT
% ============================================================
\section*{Acknowledgment}

The authors would like to thank the creators of the DIV2K dataset for
making high-resolution benchmark images publicly available, and the
Kaggle platform for providing GPU compute resources. The authors also
acknowledge the open-source PyTorch, PyWavelets, and scikit-image
communities whose tools enabled this research.

% ============================================================
% REFERENCES
% ============================================================
\bibliographystyle{IEEEtran}
\bibliography{rf}

% ============================================================
% BIOGRAPHIES
% ============================================================
\begin{IEEEbiographynophoto}{Author One}
received the B.Sc. degree in computer science from XYZ University in 2018
and the M.Sc. degree in 2020. He is currently a Ph.D. candidate. His
research interests include multimedia security, deep learning, and
signal processing.
\end{IEEEbiographynophoto}

\begin{IEEEbiographynophoto}{Author Two}
received the Ph.D. degree in electrical engineering from ABC University
in 2015. She is currently an Assistant Professor. Her research focuses
on image processing, watermarking, and computer vision.
\end{IEEEbiographynophoto}

\begin{IEEEbiographynophoto}{Author Three}
received the Ph.D. degree from DEF University in 2010. He is a Senior
Member of IEEE and currently a Professor of Electrical Engineering.
His research interests include information security, signal processing,
and machine learning.
\end{IEEEbiographynophoto}

\end{document}